\documentclass[11pt]{article}
\usepackage[margin=1in]{geometry}
\usepackage{hyperref}
\usepackage{enumitem}
\usepackage{booktabs}
\usepackage{xcolor}
\usepackage{amsmath}
\title{ProbOS --- Month 2, Week 8\\Cross-Discipline Batch 2 + Month 2 Retrospective}
\author{Nisong Monyimba}
\date{\today}

\begin{document}
\maketitle

\section*{Week 8 goal}
Show that \texttt{ParticleFilter} (Week 5) generalises across domains
the same way \texttt{MonteCarloEngine} was shown to in Month 1 Week 4:
port the ED queue and clinical trial models to a sequential filtering
setting. Close Month 2 with a retrospective and Month 3 plan.

\section*{Day-by-day plan}

\subsection*{Day 1 -- ED queue filtering setup}
\begin{itemize}[leftmargin=*]
  \item Reuse \texttt{EDQueueModel} from
    \texttt{python/examples/week4\_ed\_queue.py} unmodified
  \item Design the inference problem: infer the TRUE (unknown) arrival
    rate $\lambda$ from an observed queue-length time series, using
    \texttt{ParticleFilter}
  \item Construct a Gaussian observation log-likelihood on queue
    length, following the same pattern already established in
    \texttt{test\_particle\_filter.py} and the \texttt{/filter}
    endpoint
\end{itemize}

\subsection*{Day 2 -- ED queue filtering validation}
\begin{itemize}[leftmargin=*]
  \item Generate synthetic queue-length observations from a KNOWN true
    $\lambda$
  \item Confirm the particle filter's posterior mean for $\lambda$
    converges toward the true value as more observations arrive
  \item \texttt{python/examples/week8\_ed\_queue\_filter.py}
\end{itemize}

\subsection*{Day 3 -- Clinical trial sequential monitoring setup}
\begin{itemize}[leftmargin=*]
  \item Reuse \texttt{ClinicalTrialModel} from
    \texttt{python/examples/week4\_clinical\_trial.py} unmodified
  \item Replace Week 4's batch nested-Monte-Carlo posterior-probability
    calculation with \texttt{ParticleFilter}, updating the posterior
    belief about treatment effect after EACH patient enrols, rather
    than only at trial's end
\end{itemize}

\subsection*{Day 4 -- Clinical trial filtering validation}
\begin{itemize}[leftmargin=*]
  \item Confirm the sequentially-updated posterior at the FINAL patient
    matches (or closely agrees with) Week 4's batch calculation on the
    same trial data -- a consistency check between the two inference
    approaches on the same model
  \item \texttt{python/examples/week8\_clinical\_trial\_filter.py}
\end{itemize}

\subsection*{Day 5 -- Test coverage}
\begin{itemize}[leftmargin=*]
  \item \texttt{python/tests/test\_week8\_examples.py}: shape
    contracts, physical/statistical invariants, and validation against
    the respective ground truth for both new filtering examples --
    matching the standard already set for Week 4's cross-discipline
    examples
\end{itemize}

\subsection*{Day 6 -- pre\_week\_audit.sh + CI}
\begin{itemize}[leftmargin=*]
  \item Run \texttt{pre\_week\_audit.sh} before considering the week
    complete
  \item Extend CI's example-running steps to include the two new
    filtering scripts
\end{itemize}

\subsection*{Day 7 -- Month 2 retrospective + Month 3 plan}
\begin{itemize}[leftmargin=*]
  \item Write the Month 2 retrospective directly into
    \texttt{docs/monthly\_plans/month2/main.tex}, alongside a Month 3
    plan in \texttt{docs/monthly\_plans/month3/main.tex}
  \item Retrospective for THIS week appended to this same document
    once the week's work is actually complete, per the established
    plan-then-retrospective process
\end{itemize}

\section*{Exit criteria}
Both filtering examples pass validation against their respective
ground truth (known true $\lambda$ for the ED queue; batch-calculation
consistency for the clinical trial), with test coverage matching
Month 1's standard -- shape contracts, physical invariants, and
validation against a closed-form or independently-computed benchmark,
not just ``runs without crashing''.


\section*{What was actually accomplished (retrospective)}

\begin{itemize}[leftmargin=*]
  \item \texttt{ParticleFilter.params}: a new public property exposing
    each particle's fixed parameter values -- genuine reusable kernel
    functionality (not a one-off workaround), needed to recover the
    posterior over PARAMETERS rather than only the posterior over
    state that \texttt{PFResult.means/stds} already provided.
  \item \texttt{week8\_ed\_queue\_filter.py}: infers the true (unknown)
    patient arrival rate $\lambda$ from a noisy, observed queue-length
    time series. Starting from a deliberately WIDE
    Uniform(5, 15) prior, the filter's posterior mean landed within
    0.1--0.2 standard deviations of the true $\lambda=10.0$ across
    multiple runs.
  \item \texttt{week8\_clinical\_trial\_filter.py}: replaces Week 4's
    batch nested-Monte-Carlo posterior calculation with sequential,
    patient-by-patient updating. The filter's final posterior
    ($P(\text{treatment better}) = 0.6786$) agreed closely with Week
    4's exact Beta-Binomial batch calculation (0.6582) on the
    identical trial data.
  \item Both examples required introducing a purpose-built,
    domain-appropriate state-transition design rather than
    mechanically reusing Week 4's forward-simulation models unchanged
    -- the ED queue case needed a documented, narrowly-scoped
    exception to project doctrine around \texttt{np.random.seed()}
    (justified and defensively scoped with state save/restore); the
    clinical trial case needed a new
    \texttt{ClinicalTrialFilterModel} reflecting that real trial data
    is fully observed (deterministic state transitions), unlike the
    ED queue's genuinely noisy observations.
  \item Two real bugs were found and fixed during development, not
    glossed over: a genuine reproducibility bug in the ED queue
    ground-truth generator (not actually seeding the global state
    \texttt{EDQueueModel} depends on internally), and an
    apples-to-oranges prior mismatch in the clinical trial validation
    (comparing results computed under two different priors, which
    produced a misleading 0.18 ``disagreement'' that was actually two
    different, individually-correct answers to two different
    problems).
  \item A real CI gap was found and closed: neither new example
    script was actually being run in CI despite this plan's own Day 6
    commitment -- both are now genuinely executed in the
    \texttt{python-tests} and \texttt{integration} jobs.
\end{itemize}

\section*{Numbers at Week 8 close}
\begin{itemize}[leftmargin=*]
  \item Python tests: 368/368 passing (350 + 18 new)
  \item C++ tests: 13/13 passing
  \item mypy strict: 0 errors (full \texttt{python/} tree)
  \item ruff: 0 warnings (full declared config)
  \item \texttt{pre\_week\_audit.sh}: 16/16 checks passed
\end{itemize}

\section*{Carried into Month 3}
\texttt{ParticleFilter} is now validated across three genuinely
different domains (an exact Kalman filter comparison in Week 5, a
known-arrival-rate recovery in hospital operations, and an exact
Beta-Binomial consistency check in clinical trials) -- a stronger
domain-agnosticism claim than Month 1's forward-simulation-only
demonstration. See \texttt{docs/monthly\_plans/month2/main.tex} for
the full Month 2 retrospective and
\texttt{docs/monthly\_plans/month3/main.tex} for what comes next.
\end{document}
